\section{Problem setup}
\label{sec:setup}

\paragraph{Prediction and residuals.}
An image $x$ has a finite set $\cP$ of pixels with valid ground-truth depth $y(p)>0$, $p\in\cP$.
A frozen monocular metric model $f$ predicts $\yhat(p)=f(x)(p)$.
We work with the log-depth residual
\begin{equation}
d(p) = \log y(p) - \log \yhat(p),
\label{eq:residual}
\end{equation}
so a symmetric interval in $d$ is multiplicative in depth.

\paragraph{Sparse metric measurements.}
A sensor returns $k$ measurements per image.
At a measured pixel $p$ it reports a depth $\ytil(p)$ and hence a residual $\dtil(p)=\log\ytil(p)-\log\yhat(p)$.
For a perfect sensor $\dtil=d$; for a real one $\dtil(p)=d(p)+e(p)$ with a sensor error $e$.
The positions follow a \emph{layout}, a random rule that may depend on $x$ and $\yhat$ but never on $y$.
We use five: \emph{uniform} (a uniform sample of $\cP$), \emph{rows} (pixels on four random image rows, a few-beam scanner), \emph{cluster} (the $k$ pixels nearest one random centre, a single observed patch), \emph{edges} (uniform among edge pixels of $\yhat$, feature-point-like) and \emph{near} (uniform among the 30\% of pixels with the smallest $\yhat$, a range-limited sensor).

\paragraph{Budget split.}
Of the $k$ measurements, $k-n$ follow the layout and form the \emph{structured anchor} set $\cS\subset\cP$.
The remaining $n$ form the \emph{probe} set $\cQ$, drawn uniformly without replacement from $\cU=\cP\setminus\cS$.
Intervals are evaluated on $\cE=\cU\setminus\cQ$, the pixels that are neither anchors nor probes.
For a pixel $p$, $\delta(p)$ is its distance to the nearest measurement, divided by the image diagonal.

\paragraph{Where random probes come from.}
Probes need a sensor that can be read at random pixels: a dense or semi-dense sensor subsampled in software (phone or tablet LiDAR, time-of-flight or structured-light cameras, as in depth upsampling~\cite{baruch2021arkitscenes}), or a steerable scanner such as MEMS-mirror~\cite{pittaluga2020mems} or adaptive LiDAR~\cite{bergman2020adaptivelidar}.
A fixed-pattern scanner allows probes only among its returns, and the guarantee then covers only those pixels: on KITTI, those with a LiDAR return, not the sky.

\paragraph{Intervals and coverage.}
A method outputs $C(p)=[\ell(p),u(p)]$ for every $p\in\cE$.
The coverage of image $i$ is the fraction of its evaluation pixels that are covered,
\begin{equation}
\Cov_i=\frac{1}{|\cE_i|}\sum_{p\in\cE_i}\1\{y(p)\in C(p)\}.
\label{eq:cov}
\end{equation}
At level $1-\alpha$ we distinguish three guarantees (\cref{sec:method} states which method gives which):
\begin{itemize}[leftmargin=*,itemsep=1pt,topsep=2pt]
\item \emph{Marginal:} $\E[\Cov_i]\ge 1-\alpha$ over images, measurements and calibration data; the exchangeable unit is the image, or the scene (room, drive, visit).
\item \emph{Per-image:} $\E[\Cov_i\mid x_i,\cS_i]\ge 1-\alpha$ for \emph{every} image, only the probes random: no image is always under-covered.
\item \emph{Per distance bin:} for bins $\cB_b$ of $\delta$ (edges $0.02,0.05,0.1,0.2$), the coverage of $\cE_i\cap\{\delta\in\cB_b\}$, averaged over a new scene, is at least $1-\alpha$ in expectation for every $b$, so far pixels are not traded for near ones.
\end{itemize}

\paragraph{What is covered.}
Every guarantee concerns the values used to calibrate: sensor readings $\ytil$ certify $\ytil$, not $y$. We always \emph{evaluate} against the independent ground truth $y$.
